\documentclass[manuscript]{aastex7}
\usepackage{amsmath}
\usepackage{amssymb}
\usepackage{hyperref}

\newcommand{\vct}[1]{\boldsymbol{#1}}
\newcommand{\Er}{E_{\rm r}}
\newcommand{\Eg}{E_{\rm g}}
\newcommand{\Fr}{\vct{F}}
\newcommand{\Prad}{\mathsf{P}}
\newcommand{\Id}{\mathsf{I}}
\newcommand{\ar}{a_{\rm r}}
\newcommand{\dt}{\Delta t}
\newcommand{\dx}{\Delta x}
\newcommand{\eps}{\epsilon_{\rm mach}}
\newcommand{\chat}{\hat{c}}
\newcommand{\code}[1]{\texttt{#1}}

\begin{document}
\title{Radiation Transport and Matter Coupling in Octo-Tiger:\\
Methods and Verification Problems}

\begin{abstract}
We describe the gray two-moment radiation solver implemented in the current
Octo-Tiger development branch and distinguish its numerical verification
problems from its radiation-hydrodynamic benchmarks. The solver uses the
Levermore M1 closure, piecewise-linear reconstruction in radiation enthalpy
and a bounded radiation velocity parameter, and a two-wave HLL flux.
Radiation transport is advanced explicitly, with the physical light speed
as the default and an optional constant reduced propagation speed.
Local interaction with matter combines implicit flux damping, lagged
velocity-dependent terms, paired gas backreaction, and a backward-Euler
thermal exchange solved through a scaled scalar quartic. The current
operator sequence is first order in time for the coupled radiation update.
The verification suite includes exact streaming solutions, a periodic
Gaussian pulse compared with a linear isotropic-moment reference, and a
Gaussian luminous sphere compared with a diffusion-limit reference. Planar
wall-driven shocks and a spherical radiating blast on a Cartesian mesh extend
the suite to coupled gas and radiation evolution. We report measured
single-resolution radiation errors and short coupled implementation checks.
These results establish the tested implementation behavior, without claiming
high-resolution convergence or a completed comparison with developed
radiative-shock and blast solutions.
\end{abstract}

\section{Scope and relation to the hydrodynamic code}
Octo-Tiger evolves compressible, self-gravitating flows on an adaptive octree
using HPX task parallelism. Its finite-volume hydrodynamics, angular-momentum
treatment, fast-multipole gravity solver, and stellar-merger applications are
described by \citet{Marcello2021} and \citet{Shiber2024}. Those algorithms
provide the mesh, gas state, and distributed communication used here. The
present paper concentrates on radiation transport and the local transfer of
energy and momentum between radiation and matter. Gravity and grid rotation
are disabled in all verification calculations reported below.

The implemented closure and its useful change of variables follow
\citet{Levermore1984} and \citet{HanawaAudit2014}. The source discretization is
based on the nonrelativistic moment formulation and local implicit updates
discussed by \citet{SkinnerOstriker2013}, with the particular approximations
and time ordering specified below. Physical light speed is the default.
The optional reduced-speed-of-light approximation (RSLA) changes the
radiation time derivatives while retaining physical light speed in
moment definitions and matter forces. There is no radiation subcycling.

\section{Radiation variables and interaction model}
\subsection{Conserved moments and units}
Throughout this paper, the gray radiation moments are the lab-frame energy
density $\Er$ and the physical radiation energy flux divided by the physical
light speed $c$:
\begin{equation}
  \Fr\equiv\frac{\vct{F}_{\rm phys}}{c},\qquad
  \vct{U}=(\Er,F_x,F_y,F_z)^{\mathsf T}.
  \label{eq:state}
\end{equation}
Every component of $\vct{U}$ has energy-density units,
${\rm erg\,cm^{-3}}$ in CGS. Radiation momentum density is $\Fr/c$.
This flux normalization follows the nonrelativistic PLUTO formulation
\citep{MelonFuksman2021}; the physical flux denoted by $\vct{F}$ in
\citet{HanawaAudit2014} and \citet{SkinnerOstriker2013} corresponds to
$c\Fr$ here. The mesh arrays, checkpoint files, and exported profiles
retain the physical flux $c\Fr$; conversion to $\Fr$ is local to the
numerical kernels.

Let $\zeta=\chat/c$ be the constant parameter
\code{radiation\_reduced\_c\_factor}, with $0<\zeta\leq1$ and default
$\zeta=1$. For a stationary Cartesian grid the moment equations are
\begin{align}
 \partial_t\Er+\chat\nabla\!\cdot\Fr &= \zeta S_E,\label{eq:energy}\\
 \partial_t\Fr+\chat\nabla\!\cdot\Prad &= \zeta\vct{S}_F.\label{eq:momentum}
\end{align}
The radiation pressure tensor $\Prad$ has energy-density units. The gas
momentum is $\vct{m}=\rho\vct{v}$, and its total energy is
$\Eg=e+|\vct{m}|^2/(2\rho)$ for the ideal gas used in the coupled tests.
The local interaction sources applied to the gas are
\begin{equation}
 \left.\partial_t\Eg\right|_{\rm int}=-S_E,
 \qquad
 \left.\partial_t\vct{m}\right|_{\rm int}=-\frac{\vct{S}_F}{c}.
 \label{eq:backreaction}
\end{equation}
Thus interaction alone conserves the weighted quantities
\begin{equation}
 \Eg+\frac{\Er}{\zeta},\qquad \vct{m}+\frac{\Fr}{\chat}.
 \label{eq:weighted}
\end{equation}
These become physical total energy and momentum only for $\zeta=1$.
The same weighting gives a combined conservation law including boundary
transport. This is the RSLA convention of \citet{SkinnerOstriker2013}:
only radiation time derivatives change, and the gas sources in
Equation~(\ref{eq:backreaction}) retain their physical strength.
For $\zeta<1$, physical gas-plus-radiation energy and momentum are not
conserved quantities of the approximate equations. In particular, $c$
in the flux normalization, the radiation constant, the $v/c$ terms,
the LTE state, and the realizability condition remains the physical
light speed.
Although the equilibrium relation $\Er=\ar T^4$ is unchanged, the final
temperature of a closed gas--radiation reservoir can depend on $\zeta$
because its conserved weighted energy differs.

\subsection{Opacities, temperature, and source approximation}
We write $\chi_a=\rho\kappa_a$ for absorption and
$\chi_t=\rho\kappa_t$ for total momentum extinction, including scattering
when present. The $\kappa$ coefficients have units ${\rm cm^2\,g^{-1}}$;
the $\chi$ coefficients have units ${\rm cm^{-1}}$. A value returned by
the grid opacity routine is already a $\chi$ and must not be multiplied by
density a second time. Thermal equilibrium is
\begin{equation}
 \Er^{\rm eq}=\ar T^4,\qquad \ar=\frac{4\sigma_{\rm SB}}{c},\qquad
 T=\frac{(\gamma-1)\mu m_{\rm H}}{\rho k_{\rm B}}e,
 \label{eq:eos}
\end{equation}
where $\mu$ is the dimensionless mean molecular weight. Composition and
the thermal gas energy are evaluated on the same local cell as the
radiation source. Thermal energy excludes kinetic energy; the existing
dual-energy tracer supplies a fallback when subtracting kinetic energy
from total energy is ill-conditioned. The benchmark gas is ideal, while
the production energy extraction also removes the cold degeneracy energy
when the white-dwarf equation of state is selected.

The particular velocity-dependent source model implemented by the local
kernel is
\begin{align}
 S_E &= c\chi_a(\ar T^4-\Er)
       +(2\chi_a-\chi_t)\vct{v}\!\cdot\Fr,\label{eq:sourceE}\\
 \vct{S}_F &= -c\chi_t\Fr+
       \chi_t(\Er\Id+\Prad)\vct{v}.\label{eq:sourceF}
\end{align}
Here $\vct{v}$ is the gas velocity. These expressions retain terms through
first order in gas $v/c$ under the adopted gray-opacity assumptions. In
particular, the velocity-dependent momentum source uses the near-equilibrium
replacement $\ar T^4\simeq\Er$ that leads to the corresponding reduced
momentum equation in \citet{SkinnerOstriker2013}. It retains the full M1
pressure tensor rather than replacing it by an isotropic tensor. It is not
the complete unequal-opacity mixed-frame four-force in arbitrary
nonequilibrium conditions. This restriction matters when interpreting the
strongly heated initial state of the blast benchmark.

\subsection{Reduced-speed validity and scope}
The static-diffusion timescale criterion is
\begin{equation}
 \chat\gg v_{\max}\max(1,\tau_{\max}),\qquad
 v_{\max}=\max(|\vct{v}|+c_{\rm eff}),
 \label{eq:rsla}
\end{equation}
where $c_{\rm eff}$ includes the gas and radiation contribution to the
acoustic speed, and $\tau_{\max}$ is an optical depth across the relevant
flow or precursor scale, not simply across one cell
\citep[Section~3.2]{SkinnerOstriker2013}. This keeps radiation propagation
and diffusion faster than gas evolution when that ordering holds for
physical $c$. A constant factor preserves the weighted invariants; changing
it during a run would require additional accounting.
RSLA slows radiation transients and does not guarantee the same solution
for a rapidly heated, radiation-dominated blast. Factor and resolution
convergence must be assessed separately. The radiation-only error table
and the full-speed coupled checks below use $\zeta=1$; additional RSLA
implementation checks are identified separately. No reduced-speed
full-duration physical-validation result is claimed here.
The four radiation-only reference problems retain
$\zeta=1$ in the results workflow so that their time-dependent reference
solutions remain unchanged.

\section{Levermore M1 closure}
\subsection{Pressure tensor and physical limits}
For $\Er>0$, define the reduced flux and its direction by
\begin{equation}
 f=\frac{|\Fr|}{\Er},\qquad
 \vct{n}=\frac{\Fr}{|\Fr|}.
\end{equation}
An admissible radiation state obeys $\Er\geq0$ and $|\Fr|\leq\Er$;
vacuum has both $\Er=0$ and $\Fr=0$. Using the notation of
\citet{MelonFuksman2021}, the implemented Eddington factor is
\begin{equation}
 \xi(f)=\frac{5-2\sqrt{4-3f^2}}{3}
 =\frac{3+4f^2}{5+2\sqrt{4-3f^2}},
 \label{eq:eddington}
\end{equation}
and the pressure tensor is
\begin{equation}
 \Prad=\Er\left[
 \frac{1-\xi}{2}\Id+
 \frac{3\xi-1}{2}\vct{n}\vct{n}^{\mathsf T}\right].
 \label{eq:closure}
\end{equation}
This is the Levermore M1 closure \citep{Levermore1984}. At isotropy it gives
$\Prad=\Er\Id/3$; at free streaming it gives
$\Prad=\Er\vct{n}\vct{n}^{\mathsf T}$. The isotropic limit is evaluated
without assigning a direction to a zero flux. M1 represents a single local
preferred direction and does not independently propagate intersecting beams
that share the same cell moments.

\subsection{Enthalpy variables}
Following the reformulation of \citet{HanawaAudit2014}, introduce the
isotropic pressure component $P$, radiation enthalpy $H$, and radiation
velocity parameter $\vct{\beta}$:
\begin{align}
 H&=\Er+P,& \vct{\beta}&=\frac{\Fr}{H},\label{eq:hb}\\
 P&=\frac{1-\beta^2}{4}H,&
 \Prad&=H\vct{\beta}\vct{\beta}^{\mathsf T}+P\Id,\label{eq:hbP}\\
 \Er&=\frac{3+\beta^2}{4}H,& \Fr&=H\vct{\beta}.
 \label{eq:hbE}
\end{align}
Here $\beta^2=|\vct{\beta}|^2$. The scalar $P$ is the isotropic part of
this decomposition and is distinct from the tensor $\Prad$; $H$ is not
the comoving radiation energy. The vector $\vct{\beta}$ parametrizes the
boost of a distribution isotropic in its radiation rest frame and is
distinct from the gas velocity divided by $c$. This radiation velocity
parametrization also appears in the relativistic PLUTO module
\citep{MelonFuksman2019}.
The inverse map, written in terms of $\vct{f}=\Fr/\Er$, is
\begin{equation}
 \frac{H}{\Er}=\frac{2+\sqrt{4-3f^2}}{3},\qquad
 \vct{\beta}=\frac{3\vct{f}}{2+\sqrt{4-3f^2}}.
 \label{eq:inverse}
\end{equation}
It remains regular at zero flux. The reconstruction variables are
$\vct{q}=(H,\beta_x,\beta_y,\beta_z)$, constrained by $H\geq0$ and
$\beta^2\leq1$.
Vacuum maps to $H=0$, $\vct{\beta}=0$.

\section{Finite-volume radiation transport}
\subsection{Reconstruction and face admissibility}
The cell averages are converted to $\vct{q}$ once for each transport step.
In each coordinate direction, every component uses a generalized minmod
slope with $\theta=1.3$:
\begin{equation}
 s_i={\rm minmod}\!\left(
 \theta(q_i-q_{i-1}),\,
 \frac{q_{i+1}-q_{i-1}}{2},\,
 \theta(q_{i+1}-q_i)\right).
 \label{eq:limiter}
\end{equation}
Minmod returns the common sign times the smallest magnitude when all
arguments have that sign, and zero otherwise. The preliminary face states
are $\vct{q}_i^\pm=\vct{q}_i\pm\vct{s}_i/2$. For nonnegative neighboring
enthalpies, this scalar limiting keeps the face enthalpies nonnegative.

Limiting the three components of $\vct{\beta}$ separately does not ensure that
the reconstructed vector remains inside the unit ball. Therefore a common
factor is applied to its half-slope $\delta\vct{\beta}$.
Let $a=|\delta\vct{\beta}|^2$, $d=|\vct{\beta}\cdot\delta\vct{\beta}|$, and
$r=\max(0,1-|\vct{\beta}|^2)$. If $a+2d>r$, replace
\begin{equation}
 \delta\vct{\beta}\ \longleftarrow
 \frac{r}{\sqrt{d^2+ar}+d}\,\delta\vct{\beta}.
 \label{eq:vectorlimit}
\end{equation}
A vanishing denominator gives a zero half-slope. This expression is the
nonnegative solution of $a\lambda^2+2d\lambda=r$ and constrains both
$\vct{\beta}\pm\lambda\delta\vct{\beta}$ simultaneously. The same canonical
face state supplies the conserved jump and pressure tensor in the Riemann
solver. Roundoff excursions outside the unit ball are projected back to its
surface; they are not interpreted as a physically superluminal flux.

\subsection{Characteristics and HLL flux}
Across a face with unit normal $\vct{e}_n$, the stationary-grid flux of
Equation~(\ref{eq:state}) is
\begin{equation}
 \vct{\mathcal{F}}_n(\vct{U})=
       (\chat F_n,\chat P_{nx},\chat P_{ny},\chat P_{nz})^{\mathsf T}.
 \label{eq:pdeflux}
\end{equation}
In the $H,\vct{\beta}$ representation the two extreme characteristic speeds
are \citep{HanawaAudit2014}
\begin{equation}
 \lambda_\pm=\chat\,
 \frac{2\beta_n\pm\sqrt{(1-\beta^2)(3-\beta^2-2\beta_n^2)}}{3-\beta^2}.
 \label{eq:eigenvalues}
\end{equation}
The repeated intermediate speed is $\chat \beta_n$. The code evaluates the
radicand as $(1-\beta^2)[3(1-\beta^2)+2\beta_\perp^2]$ to avoid cancellation
near free streaming. At isotropy the extremes are $\pm\chat/\sqrt{3}$;
for a beam they collapse to its normal propagation speed.

For reconstructed left and right states, set
$s_L=\min(\lambda_-^L,\lambda_-^R)$ and
$s_R=\max(\lambda_+^L,\lambda_+^R)$. The face flux is
\begin{equation}
 \vct{\mathcal{F}}^{\rm HLL}=
 \begin{cases}
 \vct{\mathcal{F}}^L,&s_L\geq0,\\
 \vct{\mathcal{F}}^R,&s_R\leq0,\\
 \displaystyle\frac{s_R\vct{\mathcal{F}}^L-s_L\vct{\mathcal{F}}^R
       +s_Ls_R(\vct{U}^R-\vct{U}^L)}{s_R-s_L},&s_L<0<s_R.
 \end{cases}
 \label{eq:hll}
\end{equation}
The upwind branches are evaluated first, including coincident zero-speed
waves, so no zero denominator is introduced. This is a two-wave flux;
there is no separate intermediate-wave reconstruction. Radiation uses one
left/right pair at each face, independently of the hydrodynamic
face/edge/vertex reconstruction machinery.

For a moving face with normal grid velocity $w_n$, the implementation
subtracts $w_n\vct{U}$ from Equation~(\ref{eq:pdeflux}) and $w_n$ from
the characteristic speeds. With a rotating coordinate basis it also
rotates the two transverse stored flux components through the angle
$\Omega\dt$ after transport. This local rotation preserves $|\Fr|$.
All reported tests use $w_n=\Omega=0$.

\subsection{Update, storage conversion, and timestep}
The explicit transport update is the unsplit forward-Euler divergence
\begin{equation}
 \vct{U}^{\rm tr}_i=\vct{U}_i^n-
 \frac{\dt}{\dx}\sum_{d=1}^3
 \left(\vct{\mathcal{F}}_{i+\frac12\vct{e}_d}^{\rm HLL}
       -\vct{\mathcal{F}}_{i-\frac12\vct{e}_d}^{\rm HLL}\right).
 \label{eq:fv}
\end{equation}
Because the mesh stores $c\Fr$, its flux buffers use the energy component
of Equation~(\ref{eq:pdeflux}) unchanged and multiply each of the three
$F$-flux components by physical $c$. Consequently the stored flux equation
has the required $c\chat\nabla\cdot\Prad$ term. Checkpoint conversion likewise
uses an energy-density unit for $\Er$ and that unit times a velocity unit
for $c\Fr$.

Before the update, neighboring radiation ghost data are exchanged and
coarse--fine face fluxes are synchronized. Afterward, updated leaf values
are restricted to parents and physical boundaries are refreshed. The
tests reported here use uniform refinement; testing conservation of
individual restriction and reflux operations does not establish a
distributed adaptive-mesh convergence result.

Each leaf contributes the bound
\begin{equation}
 \dt_{\rm rad}=\frac{\min(C_{\rm CFL},1/4)\dx}
 {3(\chat+w_{\max})},
 \label{eq:cfl}
\end{equation}
where $w_{\max}$ bounds the absolute grid-velocity components on the
block and its faces. The factor $1/4$ is a deliberately restrictive
implementation choice for piecewise-linear reconstruction followed by
forward Euler. The shared timestep is reduced over the hydrodynamic and
radiation bounds, any requested hard upper bound, and the remaining
evolution time. Radiation is not silently advanced by extra substeps if
this shared bound is violated. At $\zeta=1$ the physical-$c$ restriction
can dominate the cost of the long-duration shock problems. RSLA relaxes
this bound but does not relax the gas or lagged-source restrictions.

The gas update completes its three Runge--Kutta stages before radiation
transport and interaction are applied once over the same $\dt$. The
gas boundaries are then refreshed, the dual-energy state is reconciled,
and the common time advances. Thus a higher-order hydrodynamic stage
sequence does not make the radiation transport, lagged source terms, or
overall splitting higher order in time. The current radiation update is
first order in time, with piecewise-linear spatial reconstruction in
smooth regions. No global second-order temporal or asymptotic-preserving
claim is made for this sequence.

\section{Local radiation--matter interaction update}
\subsection{Flux relaxation and momentum backreaction}
The interaction step begins with the post-transport radiation state and
the gas state after the hydrodynamic update. Superscript $0$ below denotes
these inputs to the local step, not necessarily the beginning of the full
hydrodynamic timestep. Density, composition, opacities, gas velocity, and
the M1 pressure tensor are held at those input values while evaluating
the lagged terms. Define $\eta_t=\chat\chi_t\dt$. Backward-Euler damping of
$\Fr$, with the velocity contribution evaluated explicitly, gives
\begin{equation}
 \Fr^1=\frac{\Fr^0+\zeta\dt\chi_t(\Er^0\Id+\Prad^0)\vct{v}^0}
                  {1+\eta_t}.
 \label{eq:fluxrelax}
\end{equation}
This update can also be evaluated as a weighted average of $\Fr^0$ and
$(\Er^0\Id+\Prad^0)\vct{v}^0/c$, using weights
$1/(1+\eta_t)$ and $\eta_t/(1+\eta_t)$.
The gas receives the paired momentum increment,
\begin{equation}
 \delta\vct{m}=\frac{\Fr^0-\Fr^1}{\chat},\qquad
 \vct{m}^1=\vct{m}^0+\delta\vct{m}.
 \label{eq:dm}
\end{equation}
This construction conserves the weighted momentum in
Equation~(\ref{eq:weighted}) locally to floating-point accuracy,
including when flux damping is stiff. The equilibrium flux and gas
force retain the physical factors of $c$.

The gas kinetic-energy change is computed without subtracting two
large squared momenta:
\begin{equation}
 \delta K=\delta\vct{m}\cdot\vct{v}^0+
                 \frac{|\delta\vct{m}|^2}{2\rho}.
 \label{eq:dK}
\end{equation}
The lagged work contribution from Equation~(\ref{eq:sourceE}) is
\begin{equation}
 W=\dt(2\chi_a-\chi_t)\vct{v}^0\cdot\Fr^0.
\end{equation}
The thermal exchange therefore starts from
\begin{equation}
 \Er^*=\Er^0+\zeta W,\qquad e^*=e^0-W-\delta K.
 \label{eq:star}
\end{equation}
Including $\delta K$ is essential: energy transferred to ordered gas
motion must not also appear as gas thermal energy.

\subsection{Backward-Euler thermal exchange}
With density and composition fixed, write
\begin{equation}
 \alpha=\ar\left[\frac{(\gamma-1)\mu m_{\rm H}}
                         {\rho k_{\rm B}}\right]^4,
 \qquad \ar T^4=\alpha e^4,
 \qquad \eta_a=\chat\chi_a\dt.
\end{equation}
The thermal subproblem is
\begin{align}
 e^1-e^*&=-\frac{\eta_a}{\zeta}[\alpha(e^1)^4-\Er^1],\\
 \Er^1-\Er^*&=\eta_a[\alpha(e^1)^4-\Er^1].
 \label{eq:BE}
\end{align}
Its conserved weighted sum $e+\Er/\zeta$ eliminates one unknown. With
\begin{equation}
 w=\frac{\eta_a}{1+\eta_a},\qquad
 q=\frac{w}{\zeta}=\frac{c\chi_a\dt}{1+\chat\chi_a\dt},\qquad
 r=e^*+q\Er^*,
\end{equation}
the remaining equation is
\begin{equation}
 x+q\alpha x^4=r,\qquad x=e^1,
 \label{eq:quartic}
\end{equation}
and the updated radiation energy is
\begin{equation}
 \Er^1=\frac{\Er^*}{1+\eta_a}+w\alpha x^4.
 \label{eq:thermalrad}
\end{equation}
For nonnegative $r$, the left-hand side of
Equation~(\ref{eq:quartic}) is monotone for $x\geq0$, so there is one
nonnegative root. Equation~(\ref{eq:thermalrad}) is evaluated as a sum
of nonnegative contributions instead of subtracting nearly equal total
and gas energies. The zero-opacity/time case is an identity update.

Large or small temperatures can make $\alpha$ impractical to form
directly. The implementation stores $\log\alpha$ and scales the root
using
\begin{equation}
 s=\min\left(r,\left[\frac{r}{q\alpha}\right]^{1/4}\right),
 \quad x=sy,\quad
 A=\frac{s}{r},\quad B=\frac{q\alpha s^4}{r},
 \label{eq:scale}
\end{equation}
so that $Ay+By^4=1$ on $0\leq y\leq1$.
Logarithmic evaluation avoids overflow in constructing the coefficients.
Both $A$ and $B$ are at most unity, and the active bounding coefficient
is assigned exactly one to retain the bracket in the nearly linear
limit. A Newton step uses derivative $A+4By^3$; if the proposal leaves
the current bracket, bisection is used instead. The bracket is updated
from the residual sign. The normalized residual tolerance is $16\eps$,
with an 80-iteration upper limit; $r=0$ is handled directly.

Finally, the total gas energy is updated from the actual radiation
energy change,
\begin{equation}
 \Eg^1=\Eg^0+\frac{\Er^0-\Er^1}{\zeta},
 \label{eq:totalgas}
\end{equation}
and the dual-energy tracer is set consistently with $e^1$.
Together, Equations~(\ref{eq:dm}) and (\ref{eq:totalgas}) make the local
momentum and energy backreaction conserve Equation~(\ref{eq:weighted}).
The quartic residual controls the agreement of the thermal and
total-energy representations.

\subsection{Positivity checks and limits of the splitting}
Implicit treatment removes the direct damping and emission/absorption
timescale restriction from the corresponding scalar relaxation steps.
It does not make the lagged work and velocity terms unconditionally
stable. The local update explicitly rejects $\Er^*<0$ or
$e^*+q\Er^*<0$, and aborts if the scaled thermal root fails to converge.
It does not clip a failed energy update, because clipping would destroy
the conservation relations. Such a failure requires a smaller shared
timestep or a change to the source discretization; the code does not
automatically retry the step at a reduced timestep.

Debug builds additionally check state positivity, vacuum consistency,
and the realizability cone, allowing a roundoff tolerance of
$4096\eps$. Those debug assertions are not all active in optimized
builds, so the application tests independently inspect the saved states.
Bounded reconstruction and implicit damping are not, by themselves, a
proof of positivity or realizability for every coupled nonequilibrium
state under the transport CFL. Accuracy of the first-order source
splitting and of the adopted $v/c$ approximation must be assessed in
the physical regime of an application.

\section{Radiation-only verification problems}
\subsection{Common setup and error definitions}
The first four tests suppress hydrodynamics and gravity. Gas variables
are held fixed rather than allowed to recoil or heat. These reference
tests use $\zeta=1$. For the Gaussian
pulse and luminous sphere the prescribed medium damps flux as
$\partial_t\Fr=-c\chi\Fr$, using
$\Fr^1=\Fr^0/(1+c\chi\dt)$ after transport. The sphere additionally
adds its prescribed physical emissivity to the energy equation.
Periodic boundaries are used for the wave, front, and Gaussian pulse;
the sphere has nonperiodic boundaries held to its reference state.
Periodic mesh communication is implemented for both gas and radiation,
although the tests in this section evolve only radiation.

Let a length unit $\ell_0$, energy-density unit $E_0$, and time unit
$\ell_0/c$ define dimensionless coordinates $\tilde{\vct{x}}$ and
$\tilde t$. The standard domain is $[-1,1]^3$ in these coordinates.
The common parameters are $\tilde\chi=5$, Gaussian width
$\tilde w=0.2$, background $\tilde E_b=1$, perturbation amplitude
$\tilde A=0.01$, and dimensionless luminosity
$\tilde L=0.01$ where appropriate. The streaming wave instead uses a
fixed relative cosine amplitude of $1/2$.

Errors are measured against cell averages of the stated reference,
not point values at cell centers. For a field $u$ and total volume $V$,
\begin{align}
 L_1(u)&=V^{-1}\sum_i V_i|u_i-u_i^{\rm ref}|,\\
 L_2(u)&=\left[V^{-1}\sum_i V_i|u_i-u_i^{\rm ref}|^2\right]^{1/2},\\
 L_\infty(u)&=\max_i|u_i-u_i^{\rm ref}|.
 \label{eq:norms}
\end{align}
These are absolute field errors and retain the units of $u$. A
dimensionless or relative error requires a separately stated scale.
For the last two tests, the reference describes an isotropic limit and
the discrepancy can contain both discretization error and closure-model
error.

\subsection{Oblique free-streaming wave}
For the integer mode vector $\vct{m}=(2,-3,1)$, domain width $L$, and
$\vct{n}=\vct{m}/\sqrt{14}$, the exact M1 solution is
\begin{equation}
 \Er(\vct{x},t)=E_0\left[1+\frac12
 \cos\!\left(\frac{2\pi}{L}
 [\vct{m}\cdot\vct{x}-c\sqrt{14}\,t]\right)\right],
 \qquad \Fr=\Er\vct{n}.
 \label{eq:wave}
\end{equation}
The integer components make the solution periodic on every Cartesian
face. Each cell-average cosine amplitude is multiplied by
$\prod_d{\rm sinc}(\pi m_d\dx/L)$, with
${\rm sinc}(z)=\sin z/z$. This test probes multidirectional transport,
the streaming closure, and conservation of all four moments.

\subsection{Periodic free-streaming front}
The initial profile is $\Er=E_0$ on the half-domain $-L/2\leq x<0$
and $\Er=10^{-10}E_0$ on $0\leq x<L/2$, repeated periodically.
The exact solution is its translation by $ct$ with
\begin{equation}
 F_x=\Er,\qquad F_y=F_z=0.
\end{equation}
The current initializer uses the streaming boundary $f=1$ exactly.
The finite-volume reference integrates the translated square wave
between cell edges, including fractional cells at either periodic
discontinuity. The test measures front broadening and phase transport.
The short results reported below stop at $\tilde t=0.2$, not after a
complete crossing of the width-$2$ domain.

\subsection{Periodic Gaussian pulse and Fourier reference}
The pulse begins with zero flux and a small periodized energy excess,
\begin{equation}
 \Er(\vct{x},0)=E_b+A
 \sum_{\vct{j}\in\mathbb{Z}^3}
 \exp\!\left[-\frac{|\vct{x}+L\vct{j}|^2}{w^2}\right],
 \qquad \Fr(\vct{x},0)=0.
\end{equation}
For the reference calculation, isotropic pressure
$\Prad=\Er\Id/3$ and prescribed scattering give
\begin{equation}
 \partial_t^2\Er+c\chi\partial_t\Er
       -\frac{c^2}{3}\nabla^2\Er=0.
 \label{eq:telegraph}
\end{equation}
Retaining the second time derivative accounts for the finite relaxation
of the initially zero flux; the pure diffusion equation would discard
that transient. For a Fourier mode with wavevector $\vct{k}$, define
$a=c\chi/2$ and $d_k^2=a^2-c^2k^2/3$. Then
\begin{align}
 A_k(t)&=e^{-at}\left[\cosh(d_kt)
                       +a\frac{\sinh(d_kt)}{d_k}\right],\\
 B_k(t)&=e^{-at}\frac{\sinh(d_kt)}{d_k},\\
 \widehat{\delta E}_k(t)&=A_k(t)\widehat{\delta E}_k(0),\\
 \widehat{\Fr}_k(t)&=-i\frac{c}{3}\vct{k}B_k(t)
                               \widehat{\delta E}_k(0).
 \label{eq:fourier}
\end{align}
Trigonometric forms apply when $d_k^2<0$, with the continuous limit at
$d_k=0$. The implementation evaluates these expressions with separate
small-argument, overdamped, and oscillatory branches to avoid cancellation
and overflow. The zero mode is constant.

The reference generator starts from error-function cell averages of the
periodized Gaussian, evolves their discrete Fourier modes, and removes
Nyquist planes on even grids so that odd spectral derivatives retain
Hermitian symmetry. Initial and final cell averages are stored together
with the reference parameters, which must match the simulation grid and
requested time. The simulation itself evolves the nonlinear M1 closure;
the Fourier solution is an exact reference for the represented isotropic
linear system, not for nonlinear M1. The small perturbation on a positive
background keeps this comparison close to the isotropic regime.

\subsection{Gaussian luminous sphere}
The spherical test has a prescribed, time-independent emissivity
\begin{equation}
 q(r)=\frac{L_b}{\pi^{3/2}w^3}e^{-r^2/w^2}.
 \label{eq:bulbsource}
\end{equation}
In a fixed scattering medium, the diffusion-limit relations
$c\nabla\cdot\Fr=q$ and $\Fr=-\nabla\Er/(3\chi)$ give
\begin{align}
 \Er(r)&=E_b+\frac{3\chi L_b}{4\pi c}
                      \frac{\operatorname{erf}(r/w)}{r},\\
 \Fr(\vct{x})&=\frac{L_b}{4\pi c}
 \frac{\operatorname{erf}(s)-2s e^{-s^2}/\sqrt{\pi}}{r^3}
 \vct{x},\qquad s=r/w.
 \label{eq:bulb}
\end{align}
Regular series evaluate the origin. Positive four-point Gauss--Legendre
quadrature in each dimension supplies reference cell averages, while
the emissivity has exact cell integrals obtained from products of
error-function differences. The prescribed source is Gaussian throughout
the calculation; it is not a uniform-emissivity solid sphere.

The initial and boundary states are taken from
Equation~(\ref{eq:bulb}). Supplying the corresponding incoming boundary
radiation is necessary; replacing it by clipped outflow would define a
different problem. The M1 evolution includes the physical source
Equation~(\ref{eq:bulbsource}) and flux damping, but uses its own nonlinear
pressure tensor. Thus this is a stationary diffusion-limit comparison,
not an exact nonlinear M1 equilibrium. It is related in purpose to
spherical illumination tests such as those discussed by
\citet{Ensman1994}, while the Gaussian source and explicit formulas above
fully specify the implemented problem.

\section{Coupled wall shocks and spherical radiating blast}
\subsection{Planar wall-driven shocks}
The two shock setups use physical parameters from the subcritical and
supercritical tests of \citet{Ensman1994}, but realize the piston through
a planar reflecting wall. The original moving spherical shell and its
equal-mass radial mesh are not reproduced. The computational cube spans
$[-3.5\times10^{10},3.5\times10^{10}]$ cm in each coordinate.
Initially
\begin{align}
 \rho_0&=7.78\times10^{-10}\ {\rm g\,cm^{-3}},\\
 v_x&=-v_p,\qquad v_y=v_z=0,\\
 T(x)&=\left[85-75\frac{x+3.5\times10^{10}\ {\rm cm}}
                       {7\times10^{10}\ {\rm cm}}\right]\ {\rm K}.
\end{align}
The piston-frame inflow speeds are $v_p=6\times10^5\ {\rm cm\,s^{-1}}$
and $2\times10^6\ {\rm cm\,s^{-1}}$, respectively. The gas has
$\gamma=5/3$ and $\mu=0.5$. Pure absorption uses
\begin{equation}
 \kappa_a=\kappa_t=\frac{3.115\times10^{-10}}{7.78\times10^{-10}}
       \ {\rm cm^2\,g^{-1}},\qquad \chi_a=\chi_t=\kappa_a\rho.
 \label{eq:shockopacity}
\end{equation}
The extinction therefore increases with compressed density; it is not
held at its initial inverse mean free path.

Initial radiation is LTE with the first-order lab-frame advective flux,
$\Er=\ar T^4$ and $F_x=(4/3)(v_x/c)\Er$. Conserved quantities are averaged
by positive tensor-product quadrature. In particular, $\langle T^4\rangle$
is used rather than $\langle T\rangle^4$. At the left wall, gas normal
momentum and radiation normal flux reverse sign in mirrored ghost cells.
The right boundary supplies the fixed upstream state at 10 K. Transverse
boundaries copy the adjacent interior values without clipping normal
momentum or radiation flux. The gas angular momentum in a ghost cell is
reconstructed consistently from its position and momentum.

The configured comparison times are $38300$ s for the subcritical case
and $7500$ s for the supercritical case. A developed subcritical solution
should exhibit a radiative precursor whose upstream gas remains cooler
than downstream equilibrium gas. In a supercritical solution the upstream
gas approaches the downstream equilibrium temperature, while a narrow
gas-temperature excess can remain at the shock. Signed radiation flux
is a useful diagnostic and must not be replaced by its magnitude when
examining the shock structure. These are physical comparison targets;
the short runs reported here have not evolved far enough to establish
either regime, and the original spherical profiles are not exact
acceptance curves for the planar wall problem.

Analytical limits provide useful temperature scales. For a strong,
cold-upstream, nonradiative wall shock, let
$\mathcal{R}=k_{\rm B}/(\mu m_{\rm H})$. The upstream speed in the shock
frame is $u_1=(\gamma+1)v_p/2$, so the ideal-gas jump conditions give
\begin{equation}
 T_2\simeq\frac{2(\gamma-1)u_1^2}{\mathcal{R}(\gamma+1)^2}
       =\frac{(\gamma-1)v_p^2}{2\mathcal{R}}.
 \label{eq:shockscale}
\end{equation}
With the present $\gamma$ and $\mu$, this gives approximately $727$ K
and $8080$ K for the two piston speeds. These are adiabatic baselines;
finite upstream temperature, radiative preheating, and net radiative
energy loss change the downstream state. In particular, $8080$ K is
not a prescribed cooled plateau for the supercritical calculation.
In the approximately steady, gas-dominated supercritical limit,
\citet[pp.~283--285]{Ensman1994} gives
$T_-\simeq T_2$ and $T_+\simeq(3-\gamma)T_2=4T_2/3$, where $T_-$ is
the gas temperature immediately upstream, $T_+$ is the narrow shock
maximum, and $T_2$ is the actual downstream equilibrium temperature.
These approximate relations do not determine the full radiative profile.

\subsection{Spherical blast on a Cartesian mesh}
The blast has spherically symmetric initial profiles sampled on the
three-dimensional Cartesian mesh. Its box spans $\pm R_o$, with
$R_o=10^{13}$ cm, and its regularized central scale is $R_0=10^{12}$ cm.
The density is
\begin{equation}
 \rho(r)=10^{-6}
 \left[\frac{\max(r,R_0)}{R_0}\right]^{-7}
              \ {\rm g\,cm^{-3}}.
 \label{eq:blastdensity}
\end{equation}
The gas is initially stationary and uses the same $\gamma$, $\mu$, and
pure-absorption mass opacity as the wall shocks. A pre-explosion
luminosity $L_0=10^{34}\ {\rm erg\,s^{-1}}$ defines an approximate LTE
diffusion atmosphere. For $R_0\leq r\leq R_o$,
\begin{align}
 E_d(r)&=\frac{L_0}{4\pi cR_o^2}
 +\frac{3\chi_0L_0}{32\pi cR_0}
       \left[\left(\frac{R_0}{r}\right)^8
                   -\left(\frac{R_0}{R_o}\right)^8\right],\\
 \Fr(r)&=\frac{L_0}{4\pi c r^2}\hat{\vct{r}},
 \qquad \chi_0=\kappa_a\rho(R_0).
 \label{eq:blastenvelope}
\end{align}
Inside the filled core,
\begin{align}
 E_d(r)&=E_d(R_0)+\frac{3\chi_0L_0}{8\pi cR_0}
                              \left(1-\frac{r^2}{R_0^2}\right),\\
 \Fr(\vct{x})&=\frac{L_0}{4\pi c R_0^3}\vct{x}.
\end{align}
At cube corners with $r>R_o$, $E_d$ is held at its edge value while the
spherical flux decreases as $r^{-2}$. The diffusion profile alone becomes
unrealizable in part of the optically thin outer envelope. Initialization
therefore sets
\begin{equation}
 \Er=\max\left[E_d,(1+10^{-12})|\Fr|\right]
\end{equation}
before computing the background gas temperature
$T=(\Er/\ar)^{1/4}$. This preserves the prescribed radial luminosity and
keeps the radiation state within the M1 cone.

An energy $E_{\rm dep}=10^{47}$ erg is then added to gas internal energy
with a Gaussian profile,
\begin{equation}
 \delta e(\vct{x})=
 \frac{E_{\rm dep}\exp(-r^2/R_0^2)}
 {\pi^{3/2}R_0^3[\operatorname{erf}(R_o/R_0)]^3}.
 \label{eq:blastdeposit}
\end{equation}
Products of error-function differences integrate this addition exactly
over each Cartesian cell. Its sum over the box is consequently
$E_{\rm dep}$ at every uniform resolution. Radiation is not given the
deposited energy directly: the subsequent matter-coupling update transfers
thermal gas energy to radiation. The luminosity used to construct the
initial atmosphere is not maintained as an energy source during the run.
The configured final time is $22000$ s.

This is a spherical explosion, although the domain and numerical grid are
Cartesian. It remains an adaptation of \citet{Ensman1994}: the original
hollow shell is filled by a regularized core, Gaussian deposition replaces
heating of one innermost Lagrangian zone, and an approximate diffusion
atmosphere replaces a relaxed transfer solution. The streaming floor,
M1 closure, and outer Cartesian boundary also affect the solution. Neither
the original breakout times nor a nonradiating Sedov--Taylor profile are
exact benchmarks for this initial-boundary-value problem. Its deposition
scale and radiation-mediated structures must be resolved before detailed
physical conclusions are drawn.

\citet[Section~5.3]{Ensman1994} supplies numerical blast profiles and
breakout diagnostics, with no exact analytical solution for that
stratified radiating problem. The Sedov--Taylor limit instead describes
an energy-conserving point explosion into cold gas of uniform density
$\rho_0$:
\begin{equation}
 R_s(t)=\xi_{\rm ST}(\gamma)\left(\frac{E_{\rm dep}t^2}{\rho_0}\right)^{1/5},
 \qquad \frac{dR_s}{dt}=\frac{2R_s}{5t},
 \label{eq:sedovlimit}
\end{equation}
where $\xi_{\rm ST}$ is a dimensionless constant fixed by the similarity
solution, distinct from the M1 Eddington factor $\xi(f)$.
The present core scale, stratified envelope, finite deposition region,
and radiative energy losses violate the assumptions of that limit;
Equation~(\ref{eq:sedovlimit}) is not an exact reference for this test.

\section{Measured checks and interpretation}
\subsection{Single-resolution radiation errors}
The current application-generated preview includes all four
radiation-only problems at $32^3$ cells. It uses
$\ell_0=3\times10^{10}$ cm and $E_0=1\ {\rm erg\,cm^{-3}}$,
so $L=6\times10^{10}$ cm, $w=6\times10^9$ cm,
$\chi=1.6666666667\times10^{-10}\ {\rm cm^{-1}}$, and
$L_b=2.698132122\times10^{29}\ {\rm erg\,s^{-1}}$ where relevant.
The final time is $0.2001384571$ s, corresponding to
$\tilde t=0.2$ with the physical light speed
$c=2.99792458\times10^{10}\ {\rm cm\,s^{-1}}$.
Table~\ref{tab:errors} reports the measured absolute energy-density norms
from those application runs. The Gaussian and sphere rows are discrepancies
from their limiting references; they are not errors against exact nonlinear
M1 solutions.

\begin{table}[ht]
\centering
\caption{Measured $32^3$ application energy-density errors at
$t=0.2001384571$ s. All entries are in ${\rm erg\,cm^{-3}}$.
\label{tab:errors}}
\begin{tabular}{lccc}
\hline
Problem & $L_1(\Er)$ & $L_2(\Er)$ & $L_\infty(\Er)$ \\
\hline
Oblique wave & $3.681505\times10^{-2}$ & $4.167722\times10^{-2}$ & $6.610035\times10^{-2}$ \\
Streaming front & $3.401321\times10^{-2}$ & $7.495241\times10^{-2}$ & $2.206418\times10^{-1}$ \\
Gaussian pulse & $1.473252\times10^{-6}$ & $8.193512\times10^{-6}$ & $1.431485\times10^{-4}$ \\
Luminous sphere & $1.461205\times10^{-5}$ & $5.905919\times10^{-5}$ & $8.268966\times10^{-4}$ \\
\hline
\end{tabular}
\end{table}

The implementation also checks all physical flux components $cF_x$,
$cF_y$, and $cF_z$, and applies fixed regression thresholds. These
single-resolution results do not determine a convergence order. In particular, no sequence of current
$32^3$, $64^3$, and $128^3$ results is presented as evidence for an
observed spatial or temporal order.

\subsection{Coupled implementation checks and conservation}
At physical light speed, each new coupled problem has completed a four-step CPU check on a
$16^3$ mesh using two HPX workers. The checks inspect all 4096 cells in
the initial and final checkpoints for finite positive matter and
radiation states and $|\Fr|\leq\Er$. Both gas and radiation must evolve;
the wall-shock checks require compression, and all three require nonzero
thermal exchange. The exported central plane must contain each expected
cell once, agree with the final checkpoint time, and satisfy the CGS
relation $\Er=\ar T_{\rm r}^4$.

The radiation conservation ledger measures the numerical boundary flux
and the actual local source increments used by the update. The four-step
application check normalizes its residual in integrated radiation energy
$\mathcal{E}$ as
\begin{equation}
 R_E=\frac{|\mathcal{E}(t)-\mathcal{E}(0)
                 +\mathcal{B}_E(t)-\mathcal{S}_E(t)|}
 {\max(|\mathcal{E}(0)|,|\mathcal{E}(t)|,
                 |\mathcal{B}_E(t)|,|\mathcal{S}_E(t)|)},
 \label{eq:ledger}
\end{equation}
where $\mathcal{B}_E$ is the cumulative outward boundary transport and
$\mathcal{S}_E$ is the cumulative radiation source contribution.
Table~\ref{tab:smoke} gives the results. This ledger checks the radiation
discrete balance with measured source terms; it is distinct from an
independent global gas-plus-radiation conservation measurement.
At reduced speed, the same ledger still measures the actual radiation
increments; the conserved combined quantities would instead be those
in Equation~(\ref{eq:weighted}).
The website uses the same unnormalized balance but divides by the fixed
initial radiation energy $|\mathcal{E}(0)|$, so growth of the source or
boundary budget cannot reduce the displayed fractional error. Its
reported fractions therefore use a different normalization from
Table~\ref{tab:smoke}.

\begin{table}[ht]
\centering
\caption{Four-step coupled CPU implementation checks on $16^3$ cells.
\label{tab:smoke}}
\begin{tabular}{lccc}
\hline
Problem & Final time (s) & Maximum $|\Fr|/\Er$ & $R_E$ \\
\hline
Subcritical wall setup & $0.04864476388$ & $0.09560791$ & $7.51\times10^{-17}$ \\
Supercritical wall setup & $0.04864476388$ & $0.09554431$ & $1.29\times10^{-16}$ \\
Spherical blast & $13.89850397$ & $0.97602206$ & $2.06\times10^{-16}$ \\
\hline
\end{tabular}
\end{table}

These short integrations test initialization, boundaries, transport,
stiff thermal exchange, checkpoint output, and reporting together.
They do not reach the configured long-duration shock or blast comparison
times. In particular, the terms subcritical and supercritical identify
the intended benchmark setups, not a demonstrated classification of
their early-time profiles.

Additional four-step CPU runs on $16^3$ cells used $\zeta=0.03$ for the
subcritical setup and $\zeta=0.1$ for the supercritical and blast setups.
They reached $1.62149213$, $0.48644764$, and $138.98503967$ s, respectively.
All passed the same full-volume state, compression where applicable,
thermal-exchange, slice, and radiation-ledger checks, and both their
initial and final checkpoints retained the requested factor. The largest
radiation-energy ledger residual with the normalization of
Equation~(\ref{eq:ledger}) was $6.43\times10^{-17}$.
The blast factor $0.1$ was used only to exercise the implementation;
it is not an accepted approximation for its physical evolution.
These runs do not test full-duration RSLA accuracy or convergence with
light-speed factor, and the physical-speed blast remains the baseline.

Additional standalone checks exercise the closure over multiple scalar
types, face realizability, transport of uniform and oblique states,
physical-flux unit conversion, local energy and momentum conservation,
opacity scaling, and the actual wall/inflow/transverse boundary assignments.
The blast initializer is checked for rotational symmetry of point states,
realizability through its outer envelope, and the deposited-energy budget
at 8, 16, and 32 cells per side. Boundary and grid checks include undefined
behavior sanitizer builds. A separate application test compares
checkpoint/restarted streaming-wave evolution with uninterrupted evolution
at $\zeta=1$ and $0.1$. Across all four radiation fields and 4096 cells,
the largest relative $L_\infty$ discrepancy was $3.15\times10^{-16}$.
The checkpoint factor overrides a conflicting input value; an actual
version-1 checkpoint correctly restores the legacy factor $1$.
Application option checks also rejected zero, negative, greater-than-one,
and nonfinite factors from both the command line and configuration file,
and verified agreement of default and explicit factor-$1$ runs.
These are continuation and option checks, not an RSLA convergence study.
GPU and multi-locality execution are outside the runtime evidence
reported here.
The RSLA update also passed all 14 selected local CTest cases, including
independent moment/source checks and production-grid transport, storage,
and interaction checks at factors $1$, $0.1$, and $0.01$ in debug,
optimized, and undefined-behavior-sanitizer configurations.

\section{Reproducibility and present limitations}
This description follows the local \code{hpx2rc-radiation} development
branch, based on \code{7485b798}, with the restored radiation module and
the subsequent local benchmark additions. The core transport and
interaction formulas are implemented in
\path{octotiger/radiation/m1.hpp}; mesh storage, flux exchange, and the
operator sequence are implemented in
\path{src/radiation/rad_grid.cpp} and the node stepping routines.
Problem definitions and cell averages are in
\path{test_problems/radiation/}, with coupled benchmark parameters in
\path{hydro_benchmarks.hpp}. The executable tests and their invocation
instructions are recorded in \path{tests/README.md}. The local results
pipeline retains the executable hash, configuration, actual final time,
physical and transport light speeds, reduced-speed factor, reference
parameters, and the distinction between application output
and results-tool fixtures.

The principal limits are the first-order radiation time sequence, the
near-equilibrium approximation in the velocity-dependent momentum
source, and the use of limiting-model references for the Gaussian pulse
and luminous sphere. Physical light-speed transport also makes some
radiation-hydrodynamic benchmarks expensive to evolve to their published
comparison times. The optional RSLA has its own timescale and convergence
requirements and does not establish those benchmarks by itself.
A completed implementation check must therefore be
kept separate from a resolved physical validation. The present results
support the numerical behaviors explicitly tested above; they do not
establish developed shock profiles, blast breakout timing, or a
high-resolution convergence rate for the coupled solver.

\appendix
\section{Archived supplemental movie links}
The following streaming-front movie links were supplied with the earlier
draft and are retained as archived supplemental material. Their
accessibility, code revision, normalization, and correspondence to the
current solver have not been reverified. They are not evidence for the
current-resolution or convergence results in this paper:
\href{https://lsu.box.com/s/o5fwovq8lcowoq0zlh5wwvkmqvq7n3ny}{$32^3$ movie},
\href{https://lsu.box.com/s/3lxru1v0uazeq2tiyfiv59c9m34y2fpz}{$64^3$ movie}, and
\href{https://lsu.box.com/s/cm7pqeqgr5dlzsr1s5tof6vvy1bw2ug3}{$128^3$ movie}.

\begin{thebibliography}{}
\bibitem[Ensman(1994)]{Ensman1994}
Ensman, L. 1994, ApJ, 424, 275--291,
\emph{Test Problems for Radiation and Radiation-Hydrodynamics Codes},
\href{https://doi.org/10.1086/173889}{doi:10.1086/173889}.

\bibitem[Hanawa \& Audit(2014)]{HanawaAudit2014}
Hanawa, T., \& Audit, E. 2014, Journal of Quantitative Spectroscopy and
Radiative Transfer, 145, 9--16,
\emph{Reformulation of the M1 Model of Radiative Transfer},
\href{https://doi.org/10.1016/j.jqsrt.2014.04.014}{doi:10.1016/j.jqsrt.2014.04.014}.

\bibitem[Levermore(1984)]{Levermore1984}
Levermore, C. D. 1984, Journal of Quantitative Spectroscopy and Radiative
Transfer, 31, 149--160,
\emph{Relating Eddington Factors to Flux Limiters},
\href{https://doi.org/10.1016/0022-4073(84)90112-2}{doi:10.1016/0022-4073(84)90112-2}.

\bibitem[Marcello et al.(2021)]{Marcello2021}
Marcello, D. C., Shiber, S., De Marco, O., et al. 2021, MNRAS, 504,
5345--5382,
\emph{Octo-Tiger: A New, 3D Hydrodynamic Code for Stellar Mergers that Uses
HPX Parallelization},
\href{https://doi.org/10.1093/mnras/stab937}{doi:10.1093/mnras/stab937}.

\bibitem[Melon Fuksman \& Mignone(2019)]{MelonFuksman2019}
Melon Fuksman, J. D., \& Mignone, A. 2019, ApJS, 242, 20,
\emph{A Radiative Transfer Module for Relativistic Magnetohydrodynamics
in the PLUTO Code},
\href{https://doi.org/10.3847/1538-4365/ab18ff}{doi:10.3847/1538-4365/ab18ff}.

\bibitem[Melon Fuksman et al.(2021)]{MelonFuksman2021}
Melon Fuksman, J. D., Klahr, H., Flock, M., \& Mignone, A. 2021,
ApJ, 906, 78,
\emph{A Two-moment Radiation Hydrodynamics Scheme Applicable to
Simulations of Planet Formation in Circumstellar Disks},
\href{https://doi.org/10.3847/1538-4357/abc879}{doi:10.3847/1538-4357/abc879}.

\bibitem[Shiber et al.(2024)]{Shiber2024}
Shiber, S., De Marco, O., Motl, P. M., et al. 2024, MNRAS, 535,
1914--1943,
\emph{Hydrodynamic Simulations of White Dwarf--White Dwarf Mergers and
the Origin of R Coronae Borealis Stars},
\href{https://doi.org/10.1093/mnras/stae2343}{doi:10.1093/mnras/stae2343}.

\bibitem[Skinner \& Ostriker(2013)]{SkinnerOstriker2013}
Skinner, M. A., \& Ostriker, E. C. 2013, ApJS, 206, 21,
\emph{A Two-moment Radiation Hydrodynamics Module in Athena Using a
Time-explicit Godunov Method},
\href{https://doi.org/10.1088/0067-0049/206/2/21}{doi:10.1088/0067-0049/206/2/21}.
\end{thebibliography}
\end{document}
